\documentclass[12pt,a4paper]{article}
\usepackage[utf8]{inputenc}
\usepackage[T1]{fontenc}
\usepackage[croatian]{babel}
\usepackage{geometry}
\usepackage{graphicx}
\usepackage{setspace}
\usepackage{titlesec}
\usepackage{fancyhdr}
\usepackage{caption}
\usepackage{longtable}
\usepackage{array}
\usepackage{tabularx}
\usepackage{amsmath}
\usepackage{amssymb}
% Times New Roman. NAPOMENA: izvorni predlozak koristi mathptmx, cije T1 varijante
% nemaju hrvatske dijakritike (c, c, s, z ispadnu iz PDF-a). newtx je Times klon s
% punim T1 pokrivanjem i vizualno je istovjetan.
\usepackage{newtxtext,newtxmath}
\usepackage{chngcntr}       % Za numeriranje po poglavljima
\usepackage{booktabs}
\usepackage{listings}
\usepackage{xcolor}
\usepackage{tikz}
\usepackage{url}
\urlstyle{tt}
\usetikzlibrary{positioning,arrows.meta,fit,calc,shapes.geometric}

% Postavke margina
\geometry{left=2.5cm,right=2.5cm,top=2.5cm,bottom=2.5cm}

% Prored 1.5
\onehalfspacing

% Bez rastavljanja rijeci na kraju retka
\hyphenpenalty=10000
\exhyphenpenalty=10000
\tolerance=9999
\emergencystretch=3em

% ---- STILOVI NASLOVA ----
\titleformat{\section}
  {\normalfont\large\bfseries\uppercase}{\thesection.}{0.5em}{}
\titleformat{\subsection}
  {\normalfont\normalsize\bfseries}{\thesubsection.}{0.5em}{}
\titleformat{\subsubsection}
  {\normalfont\normalsize\bfseries\itshape}{\thesubsubsection.}{0.5em}{}
\titleformat{\paragraph}[hang]
  {\normalfont\normalsize\itshape}{\theparagraph.}{0.5em}{}
\titlespacing*{\paragraph}{0pt}{3.25ex plus 1ex minus .2ex}{0.5em}

% Svaki \section počinje na novoj stranici
\newcommand{\sectionbreak}{\clearpage}

% ---- NUMERIRANJE SLIKA, TABLICA I JEDNADŽBI PO POGLAVLJIMA ----
\counterwithin{figure}{section}
\counterwithin{table}{section}
\counterwithin{equation}{section}

\setcounter{secnumdepth}{4}
\setcounter{tocdepth}{4}

% ---- NAZIVI AUTOMATSKIH POPISA ----
\addto\captionscroatian{%
  \renewcommand{\contentsname}{SADRŽAJ}%
  \renewcommand{\listfigurename}{POPIS SLIKA}%
  \renewcommand{\listtablename}{POPIS TABLICA}%
  \renewcommand{\refname}{LITERATURA}%
}

% ---- POTPISI SLIKA I TABLICA ----
\captionsetup[figure]{font={small,bf},labelsep=period}
\captionsetup[table]{font={small,bf},labelsep=period,position=above}

% ---- ISPISI IZ LOGA ----
\lstdefinestyle{log}{
  basicstyle=\ttfamily\scriptsize,
  breaklines=true,
  frame=single,
  rulecolor=\color{gray!50},
  backgroundcolor=\color{gray!5},
  columns=fullflexible,
  keepspaces=true
}

% ---- ZAGLAVLJE I PODNOŽJE ----
\pagestyle{fancy}
\fancyhf{}
\fancyhead[L]{\small\textit{Krešimir Hartl}}
\fancyhead[R]{\small\textit{Seminar}}
\fancyfoot[L]{\small\textit{Fakultet strojarstva i brodogradnje}}
\fancyfoot[R]{\small\textit{\thepage}}
\renewcommand{\headrulewidth}{0.4pt}
\renewcommand{\footrulewidth}{0.4pt}

\graphicspath{{slike/}}

% Plutajuci objekti smiju zauzeti vise stranice prije nego odu na svoju
\renewcommand{\topfraction}{0.9}
\renewcommand{\bottomfraction}{0.8}
\renewcommand{\textfraction}{0.07}
\renewcommand{\floatpagefraction}{0.75}
\setcounter{topnumber}{3}
\setcounter{bottomnumber}{2}
\setcounter{totalnumber}{5}

\begin{document}

% =====================================================================
% NASLOVNICA
% =====================================================================
\begin{titlepage}
    \thispagestyle{empty}
    \centering
    {\fontsize{16pt}{19pt}\selectfont SVEUČILIŠTE U ZAGREBU\\}
    {\fontsize{16pt}{19pt}\selectfont FAKULTET STROJARSTVA I BRODOGRADNJE\\}
    \vfill
    {\fontsize{28pt}{34pt}\selectfont\bfseries PROJEKTIRANJE AUTONOMNIH SUSTAVA\\}
    \vspace{0.5cm}
    {\fontsize{28pt}{34pt}\selectfont\bfseries SEMINAR\\}
    \vfill
    \noindent
    \begin{flushright}
        {\fontsize{14pt}{17pt}\selectfont\bfseries Krešimir Hartl\\}
    \end{flushright}
    \vspace{1cm}
    \begin{center}
        {\fontsize{14pt}{17pt}\selectfont Zagreb, 2026.}
    \end{center}
\end{titlepage}

% =====================================================================
% SADRŽAJ I POPISI (rimske stranice)
% =====================================================================
\pagenumbering{Roman}

\tableofcontents
\newpage
\listoffigures
\newpage
\listoftables
\newpage

% =====================================================================
% SADRŽAJ RADA (arapske stranice)
% =====================================================================
\pagenumbering{arabic}

% =====================================================================
\section{UVOD}
% =====================================================================

Cilj rada je izraditi simulacijski model dvorukog mobilnog robota i na njemu izvesti
autonomnu misiju. Model obuhvaća mobilnu bazu s dvjema robotskim rukama na vertikalnim
linearnim vodilicama i pan-tilt kamerom na vrhu, prikazan na slici~\ref{fig:sustav}, sa
senzorima, aktuatorskim mehanizmima i upravljačkim sustavima, uz parametre usklađene sa
specifikacijama stvarnog robota.

\begin{figure}[!htbp]
    \centering
    \includegraphics[width=0.44\textwidth]{s11_sustav_cad.png}
    \caption{Sustav koji je modeliran: omnidirekcijska baza, dvije vertikalne linearne vodilice s klizačima, dvije ruke sa sedam osi i pan-tilt kamera na vrhu torza}
    \label{fig:sustav}
\end{figure}

Rad je izveden u programskom okviru \textbf{ROS 2 Humble}~\cite{ros2}, u simulatoru
\textbf{Gazebo Fortress} (LTS)~\cite{gazebo}, uz \textbf{ros2\_control}~\cite{ros2control} kao
sučelje prema aktuatorima. Za planiranje gibanja ruku korišten je \textbf{MoveIt 2}~\cite{moveit},
za izgradnju karte \textbf{slam\_toolbox}~\cite{slam}, a za lokalizaciju i navigaciju
\textbf{Nav2}~\cite{nav2}.

Konačni rezultat je misija koja se pokreće \textbf{jednom naredbom}: robot se stvara u srednjoj
sobi s raširenim rukama, sam ih slaže u pozu vožnje i čeka naredbu korisnika. Nakon naredbe
odvozi se u susjednu sobu, pronalazi kutiju pomoću ArUco markera, podiže je objema rukama,
prenosi kroz vrata u treću sobu i odlaže na označeno mjesto. U posljednjem potvrđenom izvođenju
središte kutije sjelo je \textbf{5 mm} od središta odredišnog markera, uz donju plohu kutije
\textbf{4 mm} iznad ploče stola.

% =====================================================================
\section{ZADATAK I ZAHTJEVI}
\label{sec:zahtjevi}
% =====================================================================

Zadatak je razložen na \textbf{20 zahtjeva} (oznake R-01 do R-20), pri čemu svaki zahtjev ima
navedeni izvor, tehničko značenje i \textbf{mjerljiv kriterij prihvaćanja}. Zahtjevi su
grupirani u četiri skupine prikazane na slici~\ref{fig:stablo}.

\begin{figure}[!htbp]
    \centering
    \resizebox{0.88\textwidth}{!}{\input{slike/tikz/stablo_zahtjeva.tex}}
    \caption{Stablo zahtjeva: iz cilja R0 izvedene su četiri skupine i 20 pojedinačnih zahtjeva}
    \label{fig:stablo}
\end{figure}

\begin{table}[!htbp]
    \centering
    \caption{Popis zahtjeva i njihovo konačno stanje}
    \label{tab:zahtjevi}
    \scriptsize
    \begin{tabularx}{\textwidth}{@{}l X l@{}}
        \toprule
        \textbf{Oznaka} & \textbf{Zahtjev} & \textbf{Stanje} \\
        \midrule
        R-01 & Mobilna baza iz zadanog izvora                        & ispunjeno \\
        R-02 & Dvije ruke sa sedam osi i hvataljkama                 & ispunjeno \\
        R-03 & Vertikalne linearne vodilice kao aktuator             & ispunjeno \\
        R-04 & Pan-tilt mehanizam i kamera na vrhu                   & ispunjeno \\
        R-05 & Izgled robota prema priloženoj slici                  & ispunjeno \\
        R-06 & Realni parametri i realistična interakcija            & uz odstupanje \\
        R-07 & ROS 2 Humble, Gazebo Fortress, \texttt{ros2\_control} & ispunjeno \\
        R-08 & Omnidirekcijski pogon baze                            & uz odstupanje \\
        R-09 & Ruke na \texttt{ros2\_control}, upravljanje MoveIt-om & ispunjeno \\
        R-10 & Okruženje koje robot može mapirati                    & ispunjeno \\
        R-11 & Prolaz standardne širine                              & uz odstupanje \\
        R-12 & Kutija s ArUco markerom                               & ispunjeno \\
        R-13 & Zadano odredišno mjesto                               & ispunjeno \\
        R-14 & Mapiranje uz \texttt{slam\_toolbox}                   & ispunjeno \\
        R-15 & Vožnja u regiju koju zada korisnik                    & ispunjeno \\
        R-16 & Samostalan pronalazak kutije                          & ispunjeno \\
        R-17 & Podizanje kutije objema rukama                        & ispunjeno \\
        R-18 & Prolazak kroz vrata bez kutije                        & ispunjeno \\
        R-19 & Prolazak kroz vrata s kutijom                         & ispunjeno \\
        R-20 & Odlaganje kutije na zadano mjesto                     & ispunjeno \\
        \bottomrule
    \end{tabularx}
\end{table}

Od 20 zahtjeva njih 17 je ispunjeno u cijelosti, a tri uz odstupanje opisano u
poglavlju~\ref{sec:ogranicenja}.

% =====================================================================
\section{ARHITEKTURA SUSTAVA}
\label{sec:arhitektura}
% =====================================================================

\subsection{Slojevi sustava}

Sustav je izveden u šest slojeva, prikazanih na slici~\ref{fig:slojevi}. Svaki sloj prema gore
prosljeđuje \textbf{mjerenja}, a prema dolje \textbf{naredbe}, i nijedan sloj ne preskače onaj
ispod sebe. Primjer je navigacija --- misija ne šalje cilj izravno Nav2-u, nego uvijek preko
vlastitog čvora \texttt{room\_navigator}, jer samo on poznaje geometriju prolaza i zna da se u
vrata mora ući okomito.

\begin{figure}[!htbp]
    \centering
    \resizebox{0.86\textwidth}{!}{\input{slike/tikz/slojevi.tex}}
    \caption{Šest slojeva sustava; desno su komponente koje svaki sloj čine}
    \label{fig:slojevi}
\end{figure}

\subsection{Vlastiti i preuzeti paketi}

Zadatak je propisao izvore modela pojedinih komponenata, pa je dio sustava preuzet, a dio
napisan. Ta je granica u tablici~\ref{tab:paketi} povučena eksplicitno. Preuzeti paketi
\textbf{nisu} dio repozitorija --- repozitorij nosi samo manifest s točno određenim (engl.
\textit{pinned}) verzijama, a paketi se dohvaćaju izravno od autora naredbom \texttt{vcs import}.

\begin{table}[!htbp]
    \centering
    \caption{Vlastiti i preuzeti programski paketi}
    \label{tab:paketi}
    \scriptsize
    \begin{tabularx}{\textwidth}{@{}l l X@{}}
        \toprule
        \textbf{Paket} & \textbf{Podrijetlo} & \textbf{Uloga} \\
        \midrule
        \multicolumn{3}{@{}l}{\textit{Napisano za ovaj rad}} \\
        \midrule
        \texttt{pas\_dual\_arm\_bringup}        & vlastito & URDF model, svijet, sve \texttt{launch} datoteke, konfiguracije Nav2/SLAM/kontrolera, karte \\
        \texttt{pas\_dual\_arm\_scripts}        & vlastito & svi čvorovi misije, navigacije i percepcije \\
        \texttt{pas\_dual\_arm\_moveit\_config} & vlastito & MoveIt 2: SRDF, kinematika, planeri, preslikavanje kontrolera \\
        \texttt{dual\_arm\_torso}               & vlastito & opis torza i vodilica (STL model: B.~Ćaran, uz dopuštenje) \\
        \midrule
        \multicolumn{3}{@{}l}{\textit{Preuzeto, na određenim verzijama}} \\
        \midrule
        \texttt{omni\_base\_simulation} & PAL Robotics & geometrija mobilne baze, kotači, lidar \\
        \texttt{ros2\_kortex}           & Kinova       & ruke Gen3 (7 osi) i hvataljke Robotiq 2F-85 \\
        \texttt{pan\_tilt\_ros}         & I-Quotient   & pan-tilt mehanizam na vrhu torza \\
        \texttt{realsense-ros}          & Intel        & opis kamere RealSense D435 \\
        \texttt{aruco\_ros}             & PAL Robotics & preuzet, ali neiskorišten \\
        \bottomrule
    \end{tabularx}
\end{table}

Na preuzetom kodu napravljene su točno \textbf{dvije} izmjene, obje vođene kao zakrpe koje se
primjenjuju skriptom: uklanjanje triju argumenata iz opisa hvataljke kojih na korištenoj grani
nema, i dodavanje inercija te podizanje granica momenta pan-tilt zglobovima.

Paket \texttt{aruco\_ros} naveden je u zadatku kao prijedlog, ali je napisan vlastiti detektor.
Razlog je tehnički: knjižnica tog paketa ne poznaje rječnik \texttt{DICT\_4X4\_50}, koji je
odabran jer mali rječnik smanjuje vjerojatnost lažne detekcije. Vlastiti detektor koristi
\texttt{cv2.aruco} i zadržava isto sučelje prema ostatku sustava.

\subsection{Čvorovi napisani za ovaj rad}

Za ovaj je rad napisano \textbf{14 izvršnih čvorova} i sedam bibliotečnih modula, svi u jeziku
Python. Čvorovi i njihove veze prikazani su na slici~\ref{fig:graf}, a uloge su sažete u
tablici~\ref{tab:cvorovi}.

\begin{figure}[!htbp]
    \centering
    \resizebox{0.92\textwidth}{!}{\input{slike/tikz/graf_cvorova.tex}}
    \caption{Vlastiti čvorovi (bijelo) i preuzete komponente (sivo) s glavnim temama među njima}
    \label{fig:graf}
\end{figure}

\begin{table}[!htbp]
    \centering
    \caption{Čvorovi napisani za ovaj rad}
    \label{tab:cvorovi}
    \scriptsize
    \begin{tabularx}{\textwidth}{@{}l l X@{}}
        \toprule
        \textbf{Sloj} & \textbf{Čvor} & \textbf{Uloga} \\
        \midrule
        Orkestracija & \texttt{main\_task}           & stroj stanja cijele misije u osam koraka \\
        \midrule
        Navigacija   & \texttt{room\_navigator}      & vožnja soba po soba, okomit ulazak u prolaz, prekid po bočnom razmaku \\
                     & \texttt{nav\_zones}           & iz karte izvodi vrata, stolove i graf soba te gradi polje cijene \\
                     & \texttt{nav\_gui}             & upravljački panel i prikaz koraka misije \\
                     & \texttt{feature\_registry}    & prepoznavanje vrata u karti i ploča stolova u oblaku točaka \\
                     & \texttt{mapping\_tour}        & skriptirana tura kroz tri sobe pri mapiranju \\
                     & \texttt{cmd\_vel\_relay}      & jedini put naredbi brzine prema bazi i točka provedbe zaštite \\
                     & \texttt{scan\_filter}         & uklanjanje odbljesaka s vlastitog tijela iz očitanja lidara \\
                     & \texttt{footprint\_publisher} & objava stvarnog obrisa robota, uključujući nošenu kutiju \\
                     & \texttt{loc\_error}           & dijagnostika točnosti lokalizacije (bez ovlasti nad gibanjem) \\
        \midrule
        Percepcija   & \texttt{aruco\_detector}      & detekcija markera; četiri instance (glava, dva zapešća, odlaganje) \\
                     & \texttt{cloud\_restamp}       & usklađivanje koordinatnih sustava oblaka točaka za prikaz \\
        \midrule
        Poze         & \texttt{set\_posture}         & postavljanje imenovane poze uz provjeru razmaka od tla \\
                     & \texttt{table\_ready}         & priprema vodilica i ruku uz \textit{izmjerenu} potvrdu \\
        \bottomrule
    \end{tabularx}
\end{table}

\subsection{Dvije invarijante arhitekture}

\textbf{Jedan izvor očitanja okoline.} Sirovo očitanje lidara koristi se isključivo za dijagnozu.
Sve što donosi odluke --- izgradnja karte, lokalizacija, obje karte cijene, nadzor sudara i
navigacijski čvor --- čita \textbf{isti} filtrirani signal.

\textbf{Tri usklađena odgovora na pitanje širine robota.} Robot je s rukama u pozi vožnje
0,821 m širok, a prolaz je 0,980 m --- ostaje oko 8 cm po strani. Zato širina nije zapisana na
jednom mjestu, nego je izvedena na tri, koja se moraju slagati: iz geometrije modela, kao obris
koji se objavljuje kartama cijene, te kao područje koje nadzire zaštitni sloj.

% =====================================================================
\section{MODEL ROBOTA I SIMULACIJSKO OKRUŽENJE}
\label{sec:model}
% =====================================================================

\subsection{Sastavljanje modela robota}

Robot ne postoji kao gotov model --- sastavljen je iz \textbf{pet izvora} u jedan opis. Mobilna
baza, torzo s vodilicama, dvije ruke, pan-tilt mehanizam i kamera dolaze iz zasebnih paketa,
svaki sa svojim pretpostavkama o koordinatnim sustavima i mjernim jedinicama. Središnja datoteka
opisa poziva ih redom i povezuje u jedno kinematičko stablo. Isti opis robota koriste i simulator
i planer gibanja, pa oba rade nad istim modelom.

Ruke su na klizače montirane pod kutom, preko klinastih nosača, a desna je dodatno zakrenuta za
180$^\circ$ oko okomite osi. Svakoj je hvataljci dodan koordinatni sustav \texttt{tool\_tip},
pomaknut \textbf{0,127 m} od prirubnice, koji odgovara \textbf{središtu zatvorenih jastučića}.
To je važno za hvat: sve se ciljne točke zadaju upravo tom točkom, a ne prirubnicom, pa je
odstupanje koje se poslije mjeri stvarno odstupanje mjesta dodira.

\begin{figure}[!htbp]
    \centering
    \begin{minipage}{0.475\textwidth}\centering
        \includegraphics[width=\linewidth]{basic_pose.png}
    \end{minipage}\hfill
    \begin{minipage}{0.475\textwidth}\centering
        \includegraphics[width=\linewidth]{drive_pose.png}
    \end{minipage}
    \caption{Robot u simulatoru: poza pri stvaranju s raširenim rukama (lijevo) i poza vožnje \texttt{DRIVE\_V4} širine 0,821 m (desno)}
    \label{fig:poze}
\end{figure}

\begin{table}[!htbp]
    \centering
    \caption{Izmjerene dimenzije robota i mase sastavnih dijelova}
    \label{tab:mjere}
    \scriptsize
    \begin{tabularx}{\textwidth}{@{}l l X@{}}
        \toprule
        \textbf{Veličina} & \textbf{Vrijednost} & \textbf{Napomena} \\
        \midrule
        Duljina $\times$ širina $\times$ visina & 1,04 $\times$ 0,854 $\times$ 1,45 m & izmjereno nad modelom, u pozi nošenja \\
        Širina u pozi vožnje                    & 0,821 m                             & najuža provjerena poza ruku \\
        Hod vertikalnih vodilica                & 0,05 -- 0,65 m                      & prizmatični zglob, sila 1000 N \\
        Visina kamere na glavi                  & 1,39 m                              & odredila je potrebnu visinu zidova \\
        Polumjer kotača                         & 0,0762 m                            & mecanum kotači, četiri komada \\
        Masa baze                               & 34,05 kg                            & iz preuzetog modela \\
        Masa vodilice / klizača                 & 12 kg / 2 kg                        & \textbf{procjena} (v.~pogl.~\ref{sec:ogranicenja}) \\
        \bottomrule
    \end{tabularx}
\end{table}

\subsection{Senzori}

Model nosi senzore navedene u tablici~\ref{tab:senzori}. Vrijedi istaknuti dvije odluke. Prvo,
originalni laserski senzori preuzete baze su \textbf{isključeni} jer su namijenjeni starijem
simulatoru i nisu objavljivali ništa, dok su im modeli stajali u ravnini mjerenja novog senzora.
Umjesto njih ugrađen je lidar prema specifikaciji stvarnog uređaja. Drugo, kontaktni senzori na
jastučićima i senzori sile na zapešćima postoje isključivo kao \textbf{dokaz} pri ocjeni;
nijedan upravljački čvor ih ne čita, jer bi to bila privilegija koju stvarni robot nema.

\begin{table}[!htbp]
    \centering
    \caption{Senzori simulacijskog modela}
    \label{tab:senzori}
    \scriptsize
    \begin{tabularx}{\textwidth}{@{}l l X@{}}
        \toprule
        \textbf{Senzor} & \textbf{Mjesto} & \textbf{Značajke} \\
        \midrule
        Lidar                & baza      & 1080 zraka po punom krugu, doseg 0,05--25 m, razlučivost 1 mm, 25 Hz \\
        RGB-D kamera         & pan-tilt glava & 640 $\times$ 480, 15 Hz, vidno polje 1,211 rad \\
        RGB-D kamere (2)     & zapešća   & vidno polje 1,047 rad, bliska granica 0,02 m \\
        Kontaktni senzori (4)& jastučići hvataljki & 50 Hz, javljaju imena tijela u dodiru \\
        Senzori sile (2)     & zapešća   & 200 Hz, samo za vrednovanje \\
        Inercijski senzor    & baza      & 100 Hz \\
        \bottomrule
    \end{tabularx}
\end{table}

\subsection{Simulacijsko okruženje}

Svijet se sastoji od tri sobe dimenzija 6 $\times$ 6 m složene u obliku slova L, prikazane na
slici~\ref{fig:tlocrt}. Robot počinje u srednjoj (\textsc{home}) sobi, kutija se nalazi u plavoj,
a odredište u crvenoj. Takav raspored namjerno je stroži od zadatka: umjesto jednog prolaza robot
mora proći kroz \textbf{dva}, i to u oba smjera, pri čemu drugi prolaz prelazi noseći kutiju.

\begin{figure}[!htbp]
    \centering
    \resizebox{0.72\textwidth}{!}{\input{slike/tikz/tlocrt.tex}}
    \caption{Tlocrt simulacijskog svijeta s položajima stolova, kutije, markera odlaganja i
             polazišta robota}
    \label{fig:tlocrt}
\end{figure}

Zidovi su debeli 0,10 m i visoki 3,0 m --- dovoljno da kamera na visini 1,39 m ne gleda preko njih.

Prolazi su nazivne širine 1,0 m, a na izgrađenoj karti mjere \textbf{0,980 m}. Uz širinu robota od
0,821 m ostaje oko 8 cm po strani --- upravo toliko da prolaz bude stvaran test, a ne formalnost.

Kutija je kocka brida \textbf{0,30 m} i mase \textbf{0,3 kg}. Dimenzije i masa nisu bile zadane,
pa su odabrane tako da kutija bude prevelika za obuhvat jednom hvataljkom, čime se dvoručno
podizanje iz zahtjeva R-17 nameće kao jedini put. Nosi tri markera: veliki na prednjoj plohi, za
pronalaženje s udaljenosti, i dva manja na bočnim plohama, koje čitaju kamere na zapešćima
neposredno prije hvata. Bočni su markeri podignuti 0,056 m iznad središta plohe, točno u optičku
os tih kamera.

\begin{figure}[!htbp]
    \centering
    \includegraphics[width=0.44\textwidth]{box_on_desk.png}
    \caption{Kutija brida 0,30 m na stolu, s markerom za pronalaženje na prednjoj i manjim markerom za kamere zapešća na bočnoj plohi}
    \label{fig:kutija}
\end{figure}

Odredište je ploča s markerom brida 0,36 m, položena na stol u crvenoj sobi. Namjerno je šira od
kutije: kad kutija sjedne na sredinu, oko nje ostaje jednak rub sa sve četiri strane, pa se
uspješnost odlaganja vidi \textbf{golim okom}, neovisno o brojkama u zapisu.

% =====================================================================
\section{UPRAVLJANJE I PERCEPCIJA}
\label{sec:upravljanje}
% =====================================================================

\subsection{Sučelje prema aktuatorima}

Zadatak je propisao \texttt{ros2\_control} kao okvir upravljanja, što znači da se svi aktuatori
robota vode preko istog sučelja, neovisno o tome radi li se o simulaciji ili o stvarnom robotu.
Sučelje prema simulatoru izvedeno je dodatkom \texttt{ign\_ros2\_control}, a upravljačka petlja
radi na \textbf{100 Hz}.

Robotom upravlja \textbf{osam} kontrolera navedenih u tablici~\ref{tab:kontroleri}. Svi rade na
\textbf{pozicijskom} sučelju --- nigdje se ne koristi regulator sile ni momenta. To je svjesna
odluka: bez mjerenja koja bi opravdala pojačanja, uveden regulator samo bi sakrio pogreške u
geometriji iza naizgled glatkog odziva.

\begin{table}[!htbp]
    \centering
    \caption{Kontroleri i skupine zglobova kojima upravljaju}
    \label{tab:kontroleri}
    \scriptsize
    \begin{tabularx}{\textwidth}{@{}l l X@{}}
        \toprule
        \textbf{Kontroler} & \textbf{Vrsta} & \textbf{Zglobovi i napomena} \\
        \midrule
        \texttt{joint\_state\_broadcaster} & objavljivač & stanja svih zglobova \\
        \texttt{base\_controller}          & mecanum     & četiri kotača; daje odometriju i transformaciju \\
        \texttt{left\_arm\_controller}     & trajektorijski & sedam zglobova lijeve ruke \\
        \texttt{right\_arm\_controller}    & trajektorijski & sedam zglobova desne ruke \\
        \texttt{torso\_controller}         & trajektorijski & dva prizmatična klizača vodilica \\
        \texttt{pan\_tilt\_controller}     & trajektorijski & zakret i nagib glave \\
        \texttt{left\_gripper\_controller} & hvataljka   & lijeva hvataljka Robotiq 2F-85 \\
        \texttt{right\_gripper\_controller}& hvataljka   & desna hvataljka \\
        \bottomrule
    \end{tabularx}
\end{table}

\subsection{Omnidirekcijski pogon baze}

Pogon baze bio je izričit zahtjev (R-08) i ujedno prvi ozbiljan tehnički problem. Pogonski
kontroler koji koristi proizvođač baze \textbf{nije dostupan} za korištenu inačicu ROS-a 2, pa je
umjesto njega upotrijebljen \texttt{mecanum\_drive\_controller}, koji ostvaruje istu funkciju ---
neovisno gibanje po osima $x$ i $y$ te zakret --- nad istim mecanum kotačima.

Sama zamjena kontrolera nije bila dovoljna. Mecanum kotač gibanje ustranu ostvaruje \textbf{kosim
valjcima}, koji klize uzduž osi valjka, a prianjaju okomito na nju. Uz izotropno trenje kotač
takvo gibanje ne može proizvesti: baza se ili ne miče ili nekontrolirano zanosi. Rješenje je
\textbf{anizotropno trenje} --- uzdužno 0,80, poprečno 0,20, uz smjer prianjanja zadan u
koordinatnom sustavu baze. Tek je tada izmjereno čisto bočno gibanje bez zakreta, kao i vožnja
unatrag bez zanošenja.

\subsection{Prepoznavanje kutije}

Kutija se pronalazi ArUco markerima, rječnikom \texttt{DICT\_4X4\_50}. Mali rječnik odabran je
namjerno: što je manje kodova, to je veća međusobna udaljenost kodova i manja vjerojatnost lažne
detekcije pri lošem kutu ili osvjetljenju.

Detekcija je izvedena vlastitim čvorom, u \textbf{četiri instance} --- jedna na kameri glave
(marker na prednjoj plohi kutije, za pronalaženje izdaleka), po jedna na kamerama zapešća
(markeri na bočnim plohama, za posljednji korak prilaza) i jedna za marker odlaganja.

Položaj kutije ne dolazi iz jednog izvora. Uz marker, kutija se mjeri i \textbf{dubinskom
kamerom}, a dva se rezultata uspoređuju; ako se razlikuju više od 0,15 m, misija se prekida
umjesto da nastavi s nepouzdanim podatkom. Razlog za dvostruko mjerenje je praktičan: marker daje
točan položaj, ali ne i dimenzije, dok dubinsko mjerenje daje dimenzije potrebne za izračun
širine hvata.

Zanimljivo ograničenje pokazalo se pri samom prilazu. Marker se \textbf{ne} može čitati izbliza
pod velikim nagibom --- kvadrat se perspektivno sažima i detekcija otkazuje. Zato se prilaz vodi
markerom samo do približno 0,90 m, a posljednji dio puta prelazi se uz mjerenje dubinom, koje na
kut nije osjetljivo. Iz istog razloga kamere na zapešćima ne vide svoje markere u trenutku kad
jastučići već drže kutiju (na 2 cm od plohe), pa se položaj držane kutije vodi iz transformacije
zapamćene u trenutku hvatanja.

% =====================================================================
\section{MAPIRANJE I NAVIGACIJA}
\label{sec:navigacija}
% =====================================================================

\subsection{Izgradnja karte}

Karta se gradi paketom \texttt{slam\_toolbox}, uz razlučivost \textbf{0,02 m} i gušći graf poza
nego što je uobičajeno (nova poza svakih 0,1 m ili 0,1 rad). Razlog za oboje je isti: prolazi su
široki 0,98 m, a robot 0,821 m, pa pogreška karte od nekoliko centimetara nije kozmetička --- ona
odlučuje hoće li prolaz na karti uopće biti prohodan.

Mapiranje se vodi \textbf{ručno}, upravljanjem s tipkovnice, i to po pravilima koja slijede iz
trenja mecanum kotača: robot se zaustavi prije skretanja, zakrene u mjestu, pa vozi ravno ---
nikad u luku. Ruta je \textsc{home} $\rightarrow$ plava soba $\rightarrow$ \textsc{home}
$\rightarrow$ crvena soba $\rightarrow$ \textsc{home}, jer dva zatvaranja petlje drže kartu
ravnom.

\begin{figure}[!htbp]
    \centering
    \includegraphics[width=0.5\textwidth]{mapping.png}
    \caption{Izgradnja karte: laserska očitanja i zauzeće prostora tijekom ručne vožnje kroz sobe}
    \label{fig:mapiranje}
\end{figure}

\begin{figure}[!htbp]
    \centering
    \includegraphics[width=0.55\textwidth]{s61_karta.png}
    \caption{Izgrađena karta razlučivosti 0,02 m: 103,9 m$^2$ slobodne površine na rasponu
             11,9 $\times$ 11,9 m}
    \label{fig:karta}
\end{figure}

\subsection{Kada se karta smije prihvatiti}

Uobičajena mjera kvalitete karte je \textbf{pokrivenost} --- jesu li sve prostorije snimljene.
Ta mjera ovdje nije dovoljna, jer ne kaže ništa o tome jesu li prolazi na karti ostali prohodni.
Karta može pokriti sve tri sobe, a da se vrata na njoj suze za 5 cm i time potroše cijelu rezervu.

Zato je uveden \textbf{geometrijski uvjet prihvaćanja} iz tablice~\ref{tab:gate}, koji se
provjerava bez pokretanja simulacije. Ako karta padne na bilo kojem retku, ne prihvaća se nego se
tura ponavlja.

\begin{table}[!htbp]
    \centering
    \caption{Uvjet prihvaćanja karte i izmjerene vrijednosti prihvaćene karte}
    \label{tab:gate}
    \scriptsize
    \begin{tabularx}{\textwidth}{@{}X l l@{}}
        \toprule
        \textbf{Mjera} & \textbf{Prolazi ako} & \textbf{Izmjereno} \\
        \midrule
        Širina oboja vrata                    & $\ge$ 0,97 m  & \textbf{0,980 m} \\
        Odstupanje osi prolaza                & $\le$ 1 cm    & \textbf{0,0 cm} \\
        Debljina zida                         & $\le$ 0,14 m  & 0,120 m \\
        Nagib lica zida                       & $\le$ 1,0$^\circ$ & $\le$ 0,02$^\circ$ \\
        Stepenica između lica s obje strane   & $\le$ 20 mm   & $\le$ 0,9 mm \\
        \bottomrule
    \end{tabularx}
\end{table}

\subsection{Lokalizacija}

Lokalizacija je izvedena algoritmom AMCL, uz model gibanja za omnidirekcijsku bazu i od 1000 do
3000 čestica. Očekivano odstupanje mjerenja postavljeno je na \textbf{0,05 m}, u skladu sa
stvarnom točnošću ugrađenog lidara od 1 mm. Vršna pogreška lokalizacije, izmjerena usporedbom sa stvarnim položajem iz
simulatora, iznosi \textbf{3,0 cm} po osi $x$, \textbf{2,9 cm} po osi $y$ i \textbf{0,3}$^\circ$.

\subsection{Planiranje puta jedinstvenim poljem cijene}

Standardni pristup u Nav2 je \textit{sloj napuhavanja}, koji oko svake prepreke stvara područje
povećane cijene. U ovom se svijetu pokazao neupotrebljivim iz dva razloga. Prvo, upisani polumjer
robota iznosi 0,427 m, pa je u prolazu širine 0,98 m napuhavanje prekrilo \textbf{cijeli otvor}
zabranjujućom cijenom --- planer je zaključio da prolaza nema. Drugo, noge stolova lidar vidi kao
četiri odvojene prepreke, pa je oko stola nastajao uzorak nalik djetelini, a cijena se dvostruko
zbrajala.

Umjesto toga, iz karte se izravno gradi \textbf{jedinstvena ploha cijene}: zabranjujuća jezgra do
upisanog polumjera, a zatim glatka rampa širine 0,40 m s vrhom znatno ispod zabranjujuće razine.
Udaljenost se pritom mjeri od \textbf{ruba ploče stola}, a ne od nogu koje senzor vidi, pa robot
obilazi stvarnu prepreku umjesto njezine sjene. Rezultat je putanja koja se drži \textbf{0,835 m}
od stvarnih prepreka, bez ijednog umjetno zatvorenog prolaza.

\begin{figure}[!htbp]
    \centering
    \includegraphics[width=0.52\textwidth]{map.png}
    \caption{Ploha cijene izvedena iz karte: zabranjujuća jezgra oko stolova i zidova, trake nulte cijene kroz prolaze i prilazne poze (plave strelice)}
    \label{fig:polje}
\end{figure}

Dodatno su kroz prolaze i uz prilaz stolu izrezane \textbf{trake nulte cijene}. Bez njih se
potencijal globalnog planera u uskom grlu zasićuje, pa gradijentni spust gubi smjer.

\subsection{Prolazak kroz vrata}

Lokalni planer podešen je za omnidirekcijsku bazu: dopušteno je gibanje ustranu do 0,15 m/s, a
uzorkovanje brzina povećano je na 15 $\times$ 25 $\times$ 25 uzoraka po ciklusu, uz razlučivost
0,01 m/s. Uklonjena su dva dodatka koja pretpostavljaju diferencijalni pogon; jedan od njih
okretao je robota bočno za 132$^\circ$ neposredno pred vratima.

Sam prolaz, međutim, nije prepušten lokalnom planeru. Njime upravlja vlastiti čvor, koji dionicu
kroz vrata vodi kao zaseban, \textbf{okomit} ulazak: robot se prvo poravna s osi prolaza (dopušteno
odstupanje 0,08 m i 5$^\circ$), pa prolazi ravno i sporije nego po otvorenom prostoru. Tijekom
prolaza čvor neprekidno mjeri bočni razmak iz očitanja lidara i \textbf{prekida vožnju} ako
razmak padne ispod granice --- radije stane nego da okrzne dovratnik.

Prilaz stolu riješen je istom logikom. Uvedena je posebna poza ispred stola, udaljena 1,021 m od
njegova središta, odnosno 10 cm od prednjeg ruba ploče. Na toj je pozi \textbf{zabranjen zakret u
mjestu}, jer bi robot polumjera 0,673 m njime zahvatio stol; odlazak od stola uvijek počinje
vožnjom unatrag.

\begin{figure}[!htbp]
    \centering
    \includegraphics[width=0.44\textwidth]{robot_door.png}
    \caption{Robot u prolazu: objavljeni obris i laserska očitanja uz oba dovratka, pri rezervi od oko 8 cm po strani}
    \label{fig:vrata}
\end{figure}

Izmjereni rezultati: \textbf{5 od 5} pokušaja prolaska kroz vrata uspješno, uz najmanji zabilježen
bočni razmak od 3,0 do 6,1 cm pri raspoloživoj rezervi od 7,3 cm. Vožnja iz polazišta do plave
sobe traje 33,1 s, a do stola u crvenoj sobi 60,6 s.

% =====================================================================
\section{DVORUČNI HVAT I MANIPULACIJA}
\label{sec:hvat}
% =====================================================================

\subsection{Zašto hvat mora biti dvoručan}

Hvataljke Robotiq 2F-85 imaju hod prstiju \textbf{85 mm}, a brid kutije iznosi \textbf{300 mm}.
Nijedna hvataljka kutiju ne može obuhvatiti --- ni djelomično, ni za jedan brid. To nije
nedostatak izvedbe nego posljedica same postave zadatka, koji traži podizanje objema rukama.

Izabrano je rješenje u kojem obje ruke \textbf{zatvorenim} hvataljkama pritisnu suprotne bočne
plohe kutije. Jastučići zatvorenih prstiju tada služe kao dodirne plohe, a kutija se drži
pritiskom dviju ruku, ne stiskom prstiju. Takav zahvat zahtijeva da se obje ruke vode
\textbf{istovremeno}: jednostrani pritisak gura kutiju po stolu umjesto da je pridrži.

\subsection{Poze ruku}

Umjesto računanja poza iz inverzne kinematike, ključne su poze \textbf{snimljene ručno} nad
modelom, u za to napravljenom postavu, i pohranjene u registar poza. Razlog je praktičan: pozu
koja istovremeno zadovoljava doseg, izbjegavanje sudara sa stolom i prolaznost kroz vrata teže je
izračunati nego namjestiti i izmjeriti.

\begin{table}[!htbp]
    \centering
    \caption{Imenovane poze i pripadne visine klizača}
    \label{tab:poze}
    \scriptsize
    \begin{tabularx}{\textwidth}{@{}l l l X@{}}
        \toprule
        \textbf{Poza} & \textbf{Širina} & \textbf{Klizači} & \textbf{Namjena} \\
        \midrule
        \texttt{DRIVE\_V4}     & 0,821 m & 0,20 m & vožnja praznog robota kroz prolaze \\
        \texttt{DETECTION\_V4} & 1,38 m  & 0,40 m & ruke razmaknute, kamere zapešća gledaju kutiju \\
        \texttt{GRASP\_V4}     & 0,84 m  & 0,40 m & jastučići neposredno uz bočne plohe kutije \\
        \texttt{CARRY\_V4}     & 0,82 m  & 0,10 m & nošenje kutije kroz prolaz \\
        \bottomrule
    \end{tabularx}
\end{table}

Vrijedi istaknuti odnos širina iz tablice~\ref{tab:poze}. Poza hvata široka je 0,84 m, a poza
detekcije čak 1,38 m --- obje su \textbf{šire od prolaza} (0,98 m za prvu je tijesno, druga ne
prolazi uopće). Provjera nad cijelim prostorom rješenja inverzne kinematike pokazala je da je to
neizbježno: zbog položaja zapešća na osi prilaza, \textbf{svih 1620} pronađenih rješenja
vodoravnog hvata daje širinu oko 1,00 m. Zaključak je da hvat i prolaz nisu ista poza i ne mogu
to biti, pa robot nakon podizanja kutije obavezno prelazi u užu pozu nošenja.

\subsection{Slijed hvatanja}

Hvat se izvodi u koracima, pri čemu svaki korak ima vlastitu provjeru:

\begin{enumerate}\setlength\itemsep{0.1em}
    \item Klizači se dižu na 0,40 m, a ruke se \textbf{istovremeno} podižu kroz točku 12 cm iznad
          poze detekcije, čime se izbjegava zapinjanje za ploču stola.
    \item Ruke se spuštaju ravno u pozu detekcije.
    \item Baza se primiče 0,245 m, do udaljenosti s koje obje kamere na zapešćima vide svoj marker.
    \item Iz oba se markera izračuna središte kutije. Izmjereni razmak markera iznosi 0,300 m, što
          odgovara stvarnom bridu i ujedno služi kao provjera ispravnosti mjerenja.
    \item Ruke prelaze u pozu hvata, pa se vrhovi alata vode u središta markera --- ravnom
          kartezijskom linijom, s dubinom dodira od \textbf{2 mm}.
    \item Nakon dokazanog obostranog dodira uspostavlja se kruta veza, desni se jastučić rastereti
          za 2 mm, a kutija se diže \textbf{klizačima} za 0,15 m.
    \item Kutija se privuče 15 cm prema tijelu robota, ruke prelaze u pozu nošenja, baza se odmakne
          0,50 m od stola i klizači se spuštaju na 0,10 m.
\end{enumerate}

\begin{figure}[!htbp]
    \centering
    \begin{minipage}{0.475\textwidth}\centering
        \includegraphics[width=\linewidth]{pickup.png}
    \end{minipage}\hfill
    \begin{minipage}{0.475\textwidth}\centering
        \includegraphics[width=\linewidth]{pickup_vies2.png}
    \end{minipage}
    \caption{Dvoručni hvat: ruke prilaze kutiji u pozi hvata (lijevo) i jastučići naliježu na suprotne bočne plohe (desno)}
    \label{fig:hvat}
\end{figure}

Dubina dodira od 2 mm nije proizvoljna. Veći pritisak kutiju \textbf{izbacuje} --- fizikalni
mehanizam tada istovremeno rješava dodir dvaju jastučića i krutu vezu, pa objekt postaje
preodređen. Cilj dodira je dokazati kontakt, a ne stvoriti silu trenja.

Izmjerena točnost vođenja: vrh alata zaustavlja se \textbf{0,9 do 4,3 mm} od zadane točke, a
klizači dostižu zadanu visinu s pogreškom \textbf{0,0000 m}.

\subsection{Uvjet za uspostavu veze}

Kruta veza između kutije i zapešća ne uspostavlja se na temelju naredbe nego na temelju
\textbf{dokaza}: potreban je istovremen dodir na obje strane, pri čemu se broji isključivo dodir
s kutijom --- imena tijela u dodiru čitaju se iz poruke senzora, pa dodir stola ili vlastitog
tijela ne prolazi. Ako uvjet nije ispunjen, ruke se rasterete i misija se prekida.

% =====================================================================
\section{ORKESTRACIJA MISIJE}
\label{sec:misija}
% =====================================================================

\subsection{Tijek misije}

Misija je izvedena kao stroj stanja s \textbf{osam koraka}, prikazanih na slici~\ref{fig:misija}.
Robot se stvara s raširenim rukama, sam ih slaže u pozu vožnje --- namjerno \textbf{pred
korisnikom}, da se vidi da poza nije unaprijed namještena --- i zatim \textbf{čeka}. Tek nakon
naredbe korisnika kreće ostatak misije.

\begin{figure}[!htbp]
    \centering
    \resizebox{0.78\textwidth}{!}{\input{slike/tikz/misija.tex}}
    \caption{Koraci misije; svaki korak provjerava vlastiti ishod, a neuspjeh vodi u pošten
             prekid umjesto u nastavak}
    \label{fig:misija}
\end{figure}

Cijela se misija pokreće \textbf{jednom naredbom}. Sastavljanje sustava izvedeno je s vremenskim
odmakom: simulator kreće odmah, navigacijski sloj nakon 6 s, prikaz nakon 8 s, a upravljanje
zadatkom nakon 10 s.

\begin{figure}[!htbp]
    \centering
    \begin{minipage}{0.475\textwidth}\centering
        \includegraphics[width=\linewidth]{gui.png}
    \end{minipage}\hfill
    \begin{minipage}{0.475\textwidth}\centering
        \includegraphics[width=\linewidth]{robot_box_door.png}
    \end{minipage}
    \caption{Upravljački panel s prikazom koraka misije (lijevo) i prijenos kutije kroz prolaz (desno)}
    \label{fig:panel}
\end{figure}

\subsection{Odlaganje kutije}

Kad robot dođe pred odredišni stol, kutija koju nosi zaklanja marker na koji je treba
spustiti. Zato se središte markera očita ranije, s prilazne poze, i zapamti u koordinatnom
sustavu odometrije.

Ruke s kutijom zatim se podignu iznad razine ploče stola, baza se primakne stolu koliko je
moguće, a robot se bočnim gibanjem poravna sa zapamćenim središtem markera. Ruke jednim
ravnim pokretom dovedu kutiju iznad središta i spuštaju je brzinom 1 cm/s. Kad kutija sjedne
na ploču, otpušta se veza i jastučići se odmiču s njezinih ploha.

\begin{figure}[!htbp]
    \centering
    \includegraphics[width=0.46\textwidth]{drop.png}
    \caption{Kutija odložena na marker: oko nje ostaje jednak rub markera sa svih strana}
    \label{fig:odlaganje}
\end{figure}

% =====================================================================
\section{REZULTATI I ISPITIVANJA}
\label{sec:rezultati}
% =====================================================================

\subsection{Način ispitivanja}

Ispitivanje je vođeno po pravilu \textbf{jedan prirast --- jedno izvođenje}: nakon svake izmjene
pokreće se cjelovito izvođenje i njegov se ishod, uspješan ili ne, zapisuje s izmjerenim
brojevima. Tijekom rada zabilježeno je oko osamdeset takvih izvođenja.

Svaki dovršeni podsustav potvrđen je vizualno, u punom prikazu simulatora.

\subsection{Izmjereni rezultati}

\begin{table}[!htbp]
    \centering
    \caption{Ključna ispitivanja i izmjereni rezultati}
    \label{tab:rezultati}
    \scriptsize
    \begin{tabularx}{\textwidth}{@{}l X l@{}}
        \toprule
        \textbf{Ispitivanje} & \textbf{Izmjereno} & \textbf{Ishod} \\
        \midrule
        Omnidirekcijsko gibanje & čisto bočno gibanje bez zakreta, vožnja unatrag bez zanošenja & prošlo \\
        Kvaliteta karte         & vrata 0,980 m, os prolaza 0,0 cm, nagib lica $\le$ 0,02$^\circ$ & prošlo \\
        Točnost lokalizacije    & vršna pogreška 3,0 cm ($x$), 2,9 cm ($y$), 0,3$^\circ$ & prošlo \\
        Prolazak kroz vrata     & 5 od 5 prolaza; najmanji bočni razmak 3,0--6,1 cm & prošlo \\
        Vožnja do plave sobe    & 33,1 s od polazišta do ciljne poze & prošlo \\
        Vožnja do crvenog stola & 60,6 s, završna poza 4,693 m / $-$0,016 m / 3,5$^\circ$ & prošlo \\
        Točnost hvata           & vrh alata 0,9--4,3 mm od zadane točke & prošlo \\
        Dizanje klizačima       & zadano 0,5500 m, izmjereno 0,5500 m (pogreška 0,0000 m) & prošlo \\
        Odlaganje na marker     & donja ploha 4 mm iznad ploče; središte \textbf{5 mm} od markera & prošlo \\
        Negativno: kutija izvan dohvata & prekid bez uspostave veze & prošlo \\
        Negativno: marker zaklonjen      & prekid prije spuštanja, kutija ostaje u rukama & prošlo \\
        \bottomrule
    \end{tabularx}
\end{table}

\subsection{Cjelovito izvođenje misije}

Posljednje potvrđeno izvođenje provedeno je u punom prikazu simulatora. Iz zapisa su izdvojeni
redci koji nose mjerenja:

\begin{lstlisting}[style=log]
mission: drive posture ... verified
WAITING for the user: press "MISIJA: po kutiju" ...
NAV: arrived at "blue:dock"
STEP5d left tool tip 1.1 mm from its target
CARRIAGE LIFT MEASURED left=0.5500 right=0.5500 m
NAV: arrived at "red:dock"
PLACE step 6: the arms carry the cube +19.6 cm forward and +1.3 cm across
PLACE the cube bottom is +4 mm from the table top
PLACE VERIFIED: the centre of the cube is 5 mm from the marker centre
MISSION COMPLETE: the cube is on the marker.
\end{lstlisting}

Uspjeh potvrđuje tek posljednja provjera: ponovno mjerenje položaja kutije nakon što su je ruke
pustile.

\subsection{Negativna ispitivanja}

Ispitivanja u kojima sustav \textbf{mora} odustati provedena su namjerno, jer sustav koji uvijek
javlja uspjeh ne dokazuje ništa. Kutija je postavljena izvan dohvata ruku, a u drugom pokusu
zaklonjen je marker odlaganja. U oba slučaja misija je prekinuta prije nego što je učinjena
nepovratna radnja: u prvom nije uspostavljena veza s kutijom, a u drugom kutija nije puštena nego
je ostala u rukama.

% =====================================================================
\section{OGRANIČENJA I ODSTUPANJA}
\label{sec:ogranicenja}
% =====================================================================

Svako odstupanje od zadatka opisano je istim redoslijedom: što je traženo, što je pokušano, što je
pritom izmjereno, koji je uzrok i što je napravljeno umjesto toga.

\subsection{Držanje kutije trenjem}

U konačnoj izvedbi kutija se drži krutom vezom koja se uspostavlja tek nakon dokazanog obostranog
dodira. Podizanje kutije samim trenjem izvedeno je i radilo je: uz zaseban aktuatorski profil u
kojem ruke i klizači rade na sučelju sile, procijenjena sila pritiska iznosila je \textbf{3,58 N}
lijevo i \textbf{3,85 N} desno, a kutija se podizala.

Razlog zbog kojeg to rješenje nije ušlo u misiju je \textbf{nedostatak ponovljivosti}. Upravljanje
je pozicijsko, pa se sila pritiska ne zadaje nego procjenjuje, a nesigurnost te procjene prevelika
je da bi zahvat ostao stabilan kroz cijelu vožnju. Zahvat time nije stabilno određen, i to ni u
jednoj od dvije moguće izvedbe.

Pritisak \textbf{dvjema točkama}, zatvorenim jastučićima na suprotnim plohama, ne sprječava zakret
kutije. Mala razlika između ruku dovoljna je da jedna strana popusti, nakon čega kutija ispada.

Pritisak \textbf{četirima točkama}, otvorenim hvataljkama, zahtijeva istovremen dodir svih
prstiju. Kako su prsti kruti, pravac koji oni tvore mora se savršeno poravnati s plohom kutije;
svako odstupanje znači da dodir ostvaruju samo dva prsta, pa zahvat opet nije određen.

Za predaju je stoga odabrano rješenje koje je ponovljivo. Kruta veza pritom nije teleportiranje:
uspostavlja se tek kad senzori na obje strane potvrde dodir s kutijom.

\subsection{Hvataljke nisu prikladne za ovakav zadatak}

Ovo ograničenje nije posljedica izvedbe nego \textbf{izbora opreme u samom zadatku}. Hvataljka
Robotiq 2F-85 ima hod prstiju 85 mm, a zadatak traži podizanje kutije brida 300 mm. Odnos je
takav da o obuhvatu nema govora --- prsti ne mogu zahvatiti nijedan brid, niti se kutija može
uhvatiti za rub.

Posljedica je da se klasični zahvat, u kojem prsti tvore geometrijsko zatvaranje oko predmeta,
uopće ne može ostvariti. Ono što ostaje jest pritisak dviju zatvorenih šaka na suprotne plohe,
dakle zahvat koji u potpunosti ovisi o trenju i o tome da obje ruke djeluju istovremeno i
simetrično. Takav je zahvat osjetljiv na svako odstupanje: nekoliko milimetara razlike među
rukama pretvara pritisak u guranje kutije po stolu.

Za ovakav bi zadatak prikladnije bilo drukčije rješenje --- veća hvataljka ili hvataljka s
prilagodljivim prstima, vakuumski prihvat na gornjoj plohi ili podizanje podvlačenjem, poput
viličara. To, međutim, izlazi iz okvira zadanog modela robota, pa je rješenje potraženo unutar
onoga što je zadano.

\subsection{Ostala odstupanja}

\begin{table}[!htbp]
    \centering
    \caption{Preostala odstupanja od zadatka}
    \label{tab:odstupanja}
    \scriptsize
    \begin{tabularx}{\textwidth}{@{}l X X@{}}
        \toprule
        \textbf{Zahtjev} & \textbf{Odstupanje} & \textbf{Uzrok} \\
        \midrule
        R-08 & Upotrijebljen \texttt{mecanum\_drive\_controller} umjesto pogonskog kontrolera proizvođača baze & taj kontroler nije dostupan za korištenu inačicu ROS-a 2 \\
        R-11 & Prolaz je širok 0,98 m, a ne 0,80 m kako je spomenuto u zadatku & robot je s uvučenim rukama širok 0,821 m, pa je otvor od 0,80 m za njega neprohodan; umjesto toga otežana su druga mjerila --- dva prolaza, u oba smjera, sa i bez kutije \\
        R-06 & Mase vodilica i klizača su procijenjene (12 kg i 2 kg) & proizvođač vodilica ne objavljuje te podatke \\
        R-17 & Kutiju drže obje ruke, ali kruta veza postoji samo na lijevom zapešću & dvije istovremene krute veze čine kutiju preodređenom i ruše numeričko rješavanje fizike \\
        \bottomrule
    \end{tabularx}
\end{table}

\subsection{Što bi bio sljedeći korak}

Prirodan nastavak rada je uklanjanje prve stavke: ponovno uvođenje hvata trenjem, ali tek nakon
što se razriješi asimetrija između ruku, koja je i bila stvarni uzrok nestabilnosti. Uz to bi
vrijedilo zamijeniti hvataljke prihvatom primjerenim veličini predmeta, čime bi i cijeli zahvat
postao jednostavniji i robusniji. Naposljetku, mase i tromosti torza trebalo bi zamijeniti
stvarnim podacima, jer o njima izravno ovisi vjerodostojnost svake dinamičke simulacije.

% =====================================================================
\section{ZAKLJUČAK}
% =====================================================================

U radu je izrađen cjelovit simulacijski model dvorukog mobilnog robota i na njemu izvedena
autonomna misija koja se pokreće jednom naredbom: robot čeka naredbu korisnika, odvozi se u
zadanu sobu, pronalazi kutiju, podiže je objema rukama, prenosi kroz vrata i odlaže na označeno
mjesto. Središte kutije pritom sjeda \textbf{5 mm} od središta odredišnog markera.

Od 21 postavljenog zahtjeva ispunjeno je 17 u cijelosti, a tri uz odstupanja koja su opisana i
obrazložena. Tehnički su najzahtjevnija bila dva dijela. Prvi je prolazak kroz otvor koji je od
robota širi svega osam centimetara po strani, riješen jedinstvenom plohom cijene umjesto
uobičajenog napuhavanja prepreka i vođenjem prolaza kao zasebne, okomite dionice. Drugi je
dvoručni hvat predmeta koji je tri i pol puta veći od raspona hvataljke, riješen istovremenim
pritiskom dviju zatvorenih šaka uz uspostavu veze tek nakon dokazanog obostranog dodira.

Metodološki, najkorisnijim se pokazalo pravilo da se svaka tvrdnja mora moći izmjeriti. Nekoliko
problema koji su dugo izgledali kao pogreške u upravljanju --- klizači koji se ne dižu, ruka koja
prolazi kroz nogu stola, lokalizacija pomaknuta za nekoliko centimetara --- imalo je uzrok u
jednoj krivo postavljenoj vrijednosti ili u pojednostavljenom modelu okoline, a otkriveni su tek
kad su provjere prestale ovisiti o naredbi koju sustav sam izdaje. Iz istog razloga sustav radije
prekida izvođenje nego što nastavlja s nepouzdanim podatkom.

Model je time spreman kao podloga za daljnji rad: za uvođenje upravljanja silom pri hvatu, za
zamjenu prihvata primjerenijeg veličini predmeta i, konačno, za prijenos istih upravljačkih
slojeva na stvarni robot, budući da je cijeli sustav izgrađen na sučeljima koja se pri tom
prijenosu ne mijenjaju.

% =====================================================================
% LITERATURA
% =====================================================================
\newpage
\begin{thebibliography}{99}
\addcontentsline{toc}{section}{LITERATURA}

\bibitem{zadatak}
B.~Ćaran, \textit{Izrada simulacijskog modela DUAL ARM robota u Gazebo okruženju},
zadatak kolegija Projektiranje autonomnih sustava, Fakultet strojarstva i brodogradnje,
Sveučilište u Zagrebu, 30.~11.~2025.

\bibitem{mail}
B.~Ćaran, \textit{PAS -- Dual arm robot -- Podaci}, dopis s izvorima modela i opisom
scenarija, 4.~5.~2026.

\bibitem{ros2}
Open Robotics, \textit{ROS 2 Documentation: Humble Hawksbill},
\url{https://docs.ros.org/en/humble/}

\bibitem{gazebo}
Open Robotics, \textit{Gazebo Fortress Documentation},
\url{https://gazebosim.org/docs/fortress/}

\bibitem{ros2control}
\textit{ros2\_control: Control Framework for ROS 2},
\url{https://control.ros.org/humble/}

\bibitem{moveit}
PickNik Robotics, \textit{MoveIt 2 Documentation},
\url{https://moveit.picknik.ai/humble/}

\bibitem{nav2}
S.~Macenski, F.~Martín, R.~White, J.~Ginés Clavero, \textit{The Marathon 2: A Navigation System},
IEEE/RSJ International Conference on Intelligent Robots and Systems (IROS), 2020.

\bibitem{slam}
S.~Macenski, I.~Jambrecic, \textit{SLAM Toolbox: SLAM for the dynamic world},
Journal of Open Source Software, sv.~6, br.~61, 2021.

\bibitem{aruco}
S.~Garrido-Jurado, R.~Muñoz-Salinas, F.~J.~Madrid-Cuevas, M.~J.~Marín-Jiménez,
\textit{Automatic generation and detection of highly reliable fiducial markers under occlusion},
Pattern Recognition, sv.~47, br.~6, str.~2280--2292, 2014.

\bibitem{omnibase}
PAL Robotics, \textit{omni\_base\_simulation}, licenca Apache-2.0, verzija \texttt{77248ac},
\url{https://github.com/pal-robotics/omni\_base\_simulation}

\bibitem{kortex}
Kinova Robotics, \textit{ros2\_kortex}, licenca BSD, verzija \texttt{116d87a},
\url{https://github.com/Kinovarobotics/ros2\_kortex}

\bibitem{pantilt}
I-Quotient Robotics, \textit{pan\_tilt\_ros}, licenca MIT, verzija \texttt{9b08758},
\url{https://github.com/I-Quotient-Robotics/pan\_tilt\_ros}

\bibitem{realsense}
Intel RealSense, \textit{realsense-ros}, licenca Apache-2.0, verzija \texttt{6d87b07},
\url{https://github.com/realsenseai/realsense-ros}

\bibitem{arucoros}
PAL Robotics, \textit{aruco\_ros}, licenca MIT, verzija \texttt{86a0bbb},
\url{https://github.com/pal-robotics/aruco\_ros}

\bibitem{predavanja}
\textit{Materijali kolegija Projektiranje autonomnih sustava}, Fakultet strojarstva i
brodogradnje, Sveučilište u Zagrebu, 2026.

\end{thebibliography}

% =====================================================================
% PRILOZI
% =====================================================================
\newpage
\section*{PRILOZI}
\addcontentsline{toc}{section}{PRILOZI}

\subsection*{Prilog A: Repozitorij}

Cjelokupan izvorni kod, opis modela, konfiguracije, karta i upute za pokretanje nalaze se u
javnom repozitoriju:

\begin{center}
\url{https://github.com/KxHartl/PAS-DUAL-ARM}
\end{center}

Datoteka \texttt{README.md} sadrži preduvjete, postupak instalacije i pokretanja te popis
vanjskih paketa s autorima, licencama i točnim verzijama. Detaljnije upute podijeljene su u
\texttt{RUNNING.md} (rad po terminalima) i \texttt{MAPPING.md} (izgradnja karte od nule).

\subsection*{Prilog B: Pokretanje misije}

Preduvjeti su Ubuntu 22.04, ROS 2 Humble, Gazebo Fortress i pripadni paketi navedeni u
\texttt{README.md}. Nakon dohvata repozitorija:

\begin{lstlisting}[style=log]
vcs import src < ros2.repos      # dohvat vanjskih paketa na zadanim verzijama
./scripts/apply_patches.sh       # dvije dokumentirane zakrpe
rosdep install --from-paths src --ignore-src -y -r
./scripts/run_native.sh colcon build --symlink-install
\end{lstlisting}

Cijela se misija zatim pokreće jednom naredbom:

\begin{lstlisting}[style=log]
./scripts/run_native.sh ros2 launch pas_dual_arm_bringup mission.launch.py
\end{lstlisting}

Kad u zapisu piše \texttt{WAITING for the user}, na upravljačkom se panelu pritisne zeleni gumb
\textbf{,,MISIJA: po kutiju``}. Sve dalje ide samostalno, do retka
\texttt{MISSION COMPLETE}. Karta je već u repozitoriju, pa mapiranje nije potrebno za
pokretanje.

\subsection*{Prilog C: Datoteke za pokretanje}

\begin{table}[!htbp]
    \centering
    \scriptsize
    \begin{tabularx}{\textwidth}{@{}l X@{}}
        \toprule
        \textbf{Datoteka} & \textbf{Što pokreće} \\
        \midrule
        \texttt{mission.launch.py}     & cijela misija iz jedne naredbe (simulator, navigacija, prikaz, zadatak) \\
        \texttt{sim.launch.py}         & simulator, model robota, most prema ROS-u i osam kontrolera \\
        \texttt{nav2.launch.py}        & karta, lokalizacija, Nav2, zone, navigacijski čvor i panel \\
        \texttt{task.launch.py}        & MoveIt, detekcija markera i upravljanje zadatkom \\
        \texttt{mapping.launch.py}     & \texttt{slam\_toolbox} i prikaz za izgradnju karte \\
        \texttt{aruco.launch.py}       & četiri instance detektora markera \\
        \texttt{display.launch.py}     & samo prikaz modela robota, bez simulacije \\
        \bottomrule
    \end{tabularx}
\end{table}

\subsection*{Prilog D: Provjere bez simulatora}

Prije svakog pokretanja simulacije mogu se, i trebaju, pustiti provjere koje ne zahtijevaju
simulator --- one provjeravaju je li karta prihvatljiva i jesu li navigacijske zone ispravno
izvedene iz nje:

\begin{lstlisting}[style=log]
python3 scripts/check_doors.py        # vrata prepoznata iz karte
python3 scripts/check_zones.py        # zone, portalne i prilazne poze
python3 scripts/check_map_geometry.py # geometrijski uvjet prihvacanja karte
\end{lstlisting}

\end{document}
